\section{Extended Profile and Trust-Repair Results}
\label{app:extended_results}

This appendix reports final-linear gain sweeps, computational profile diagnostics, and trust-repair analyses.

\subsection{Alpha sweep}
\label{sec:phenotypes_alpha}

Sweeping $\alpha$ across $\{0.05, 0.1, 0.3, 0.5, 1.0, 2.0, 4.0, 8.0\}$ tests whether precision gain controls reactive social confidence across metrics. The average within-episode max--min range of posterior mean $\bar{\beta}_k$ increases monotonically with $\alpha$ in the betrayal regime:

\begin{table}[H]
\centering
\caption{Mean within-episode $\bar{\beta}_k$ range by precision-gain $\alpha$ (betrayal regime).}
\label{tab:alpha_sweep}
\begin{tabular}{cc}
\toprule
$\alpha$ & Mean within-episode $\bar{\beta}_k$ range \\
\midrule
0.05 & 0.097 \\
0.1  & 0.126 \\
0.3  & 0.195 \\
0.5  & 0.219 \\
1.0  & 0.277 \\
2.0  & 0.402 \\
4.0  & 0.495 \\
8.0  & 0.675 \\
\bottomrule
\end{tabular}
\end{table}

\begin{figure}[H]
\centering
\appendixfigure{fig_alpha_sweep.pdf}
\caption{Precision-gain sweep across open graded and betrayal environments. Higher $\alpha$ expands the average within-episode max--min range of posterior mean $\bar{\beta}_k$, while engagement, recovery, entropy, and payoff readouts vary non-monotonically. Precision gain therefore controls confidence-reactivity amplitude rather than mapping directly to reward.}
\label{fig:alpha_sweep_appendix}
\end{figure}
\appendixfloattosection

\subsection{Gain-prior profile signatures}
\label{sec:phenotypes_quadrants}

``Naive'' and ``cautious'' refer to priors over $q(\beta_k)$: naive $[0.40, 0.40, 0.15, 0.04, 0.01]$, cautious $[0.01, 0.04, 0.15, 0.40, 0.40]$. Low/high gain use $\alpha=0.1/3.0$. Profile labels correspond to these combinations: naive-stubborn uses naive prior/low gain, anxious-reactive uses naive prior/high gain, avoidant-rigid uses cautious prior/low gain, and hypervigilant uses cautious prior/high gain. The default uses $\alpha=3.0$ with centered initialization at $\beta_k=1.0$. These labels denote computational parameter combinations rather than clinical classifications or validated human phenotypes. Table~\ref{tab:phenotypes} reports raw metrics and Figure~\ref{fig:phenotype_quadrants_appendix} shows column-normalized profile signatures.

\begin{table}[H]
\centering
\caption{Behavioral metrics by profile (betrayal environment). Default agent shown for reference.}
\label{tab:phenotypes}
\resizebox{\textwidth}{!}{%
\begin{tabular}{lccccc}
\toprule
Profile & Payoff & Recovery rounds & Gini & Mean within-episode $\bar{\beta}_k$ range & Trust asymmetry \\
\midrule
Anxious-reactive   & 2285.8 & 25.1 & 0.626 & 0.459 & 17.72 \\
Hypervigilant      & 2230.7 & 21.5 & 0.609 & 0.556 & 7.25 \\
Naive-stubborn     & 2218.0 & 17.8 & 0.567 & 0.173 & 5.46 \\
Avoidant-rigid     & 2166.0 & 15.7 & 0.567 & 0.302 & 1.75 \\
Default            & 2244.2 & 32.9 & 0.643 & 0.427 & 8.11 \\
\bottomrule
\end{tabular}
}
\end{table}

\begin{figure}[H]
\centering
\appendixfigure{fig_phenotype_quadrants.pdf}
\caption{Gain-prior profile signatures from the prior-by-$\alpha$ factorial (betrayal environment). Cells are column-normalized across profiles for visual comparison; raw metric values are reported in Table~\ref{tab:phenotypes}. The heatmap shows that gain and prior model fitness produce distinct calibration signatures rather than a single reward-aligned ordering.}
\label{fig:phenotype_quadrants_appendix}
\end{figure}
\appendixfloattosection

\subsection{Forgiveness and trust repair}

The forgiveness protocol uses baseline rounds 1--80, betrayal rounds 81--120, and repair/reversion rounds 121--200. Table~\ref{tab:forgiveness} reports reengagement, payoff recovery, latency, and mean within-episode $\bar{\beta}_k$ range; Figure~\ref{fig:forgiveness_appendix} shows reengagement, $\beta_k$ recovery, and payoff recovery.

\Needspace{10\baselineskip}
\begin{table}[H]
\centering
\caption{Forgiveness and trust-repair metrics by precision profile. Reengagement is the post-repair P0 selection rate; payoff recovery compares late post-repair payoff with the late pre-betrayal baseline.}
\label{tab:forgiveness}
\resizebox{\textwidth}{!}{%
\begin{tabular}{lcccc}
\toprule
Variant & Reengagement & Payoff recovery & Latency & Mean within-episode $\bar{\beta}_k$ range \\
\midrule
cautious-high-$\alpha$ & 0.530 & 1.047 & 14.6 & 0.607 \\
cautious-low-$\alpha$ & 0.686 & 1.068 & 6.5 & 0.301 \\
default & 0.586 & 1.030 & 10.4 & 0.485 \\
naive-high-$\alpha$ & 0.579 & 1.000 & 3.6 & 0.604 \\
naive-low-$\alpha$ & 0.615 & 1.022 & 7.4 & 0.190 \\
no-affect & 0.587 & 1.049 & 12.0 & -- \\
\bottomrule
\end{tabular}
}
\end{table}

\begin{figure}[H]
\centering
\appendixfigure{fig_forgiveness.pdf}
\caption{Forgiveness/trust-repair diagnostic from the profile analyses. Vertical markers indicate the transition from cooperative baseline to betrayal and then to repair/reversion. Reengagement and payoff recovery are shown alongside the reverted-partner $\beta_k$ trajectory. Because $\beta_k$ is inverse precision, lower $\beta_k$ indicates higher confidence / higher policy precision; no-affect is omitted from the $\beta_k$ trajectory panel because it has no $\beta$ tracker. The figure shows that agents can reengage a repaired partner before rebuilding the same confidence dynamics.}
\label{fig:forgiveness_appendix}
\end{figure}
